\section{Corpus et données}\label{sec:corpus}

\subsection{Les PDF sources}
Le tableau~\ref{tab:corpus} liste les fichiers. Plusieurs PDF regroupent un code et des textes annexes (lois, décrets,
conventions) dont la numérotation des articles recommence à 1 : « article 12 » seul est donc ambigu dans ces fichiers.

\begin{table}[H]
\centering
\small
\begin{tabular}{lrrr}
\toprule
Fichier source & Pages & \emph{Chunks} (index) & Extraction PDF \\
\midrule
ArbitrageFr & 68 & 131 & vérifiée \\
COC & 267 & 1407 & vérifiée \\
CommerceFr & 219 & 561 & vérifiée \\
DIPFR & 19 & 78 & vérifiée \\
dtreelFr & 152 & 535 & vérifiée \\
EnfantFr & 142 & 290 & vérifiée \\
NationaliteFr & 22 & 69 & vérifiée \\
PenalFr & 100 & 360 & vérifiée \\
Procedurecivilecom & 124 & 469 & vérifiée \\
ProcedurepenaleFr & 146 & 492 & vérifiée \\
societeFr & 204 & 593 & vérifiée \\
Statutpersonnel & 148 & 408 & vérifiée \\
travailFr & 150 & 496 & vérifiée \\
tun\_RecAcFreeLi & 15 & 142 & vérifiée \\
\bottomrule
\end{tabular}
\caption{Les quatorze PDF sources, leurs pages, le nombre de \emph{chunks} de l'index et le statut de l'extraction indépendante (pypdf).}
\label{tab:corpus}
\end{table}


\subsection{Découpage par article et identifiants stables}
L'index découpe chaque PDF en \term{chunks} d'un article. Un \term{chunk} ne commence qu'à un véritable intitulé
(« Article 218 », « Art. 22 », « Article premier », avec suffixes \emph{bis}, \emph{ter}, etc.) ; un renvoi dans le texte
(« prévues à l'article 62 du présent code ») ne coupe plus l'article. Les textes regroupés sont séparés en sections
(la numérotation qui repart à 1 ouvre une section). Chaque \term{chunk} reçoit un identifiant stable de la forme
\texttt{source/section/article}, par exemple \texttt{PenalFr/s02/226-ter}. Comparer les réponses au niveau de cet
identifiant, et non du seul numéro d'article, évite de confondre l'article 12 de deux codes ou de deux textes regroupés.

\subsection{Vérification des \emph{chunks} contre les PDF}
Pour ne pas évaluer le système contre ses propres données, nous avons extrait de nouveau le texte des PDF avec un autre
outil (pypdf, alors que l'index utilise PyMuPDF) et un découpage indépendant par expression régulière. Le script
\texttt{eval/audit\_chunks.py} compare les deux côtés.

\begin{itemize}
  \item \textbf{Couverture.} \NAuditOnlypdf{} étiquettes d'articles n'apparaissent que dans le PDF et \NAuditOnlyindex{}
        que dans l'index (tableau~\ref{tab:audit}). Une partie de ces écarts vient de notre propre expression régulière
        (intitulés atypiques), une autre de l'index.
  \item \textbf{Fidélité.} Sur \NAuditN{} \term{chunks} tirés au hasard (graine \NSeed), \NAuditChecked{} ont pu être
        comparés ; \NAuditMismatch{} ne correspondent pas au texte du PDF (\NAuditRate~\%), et \NAuditNopdf{} n'ont pas
        d'article PDF correspondant (numérotation des sections différente). Les écarts observés sont des marques de notes
        de bas de page perdues par l'un des extracteurs et, dans le code de commerce, une ligne de pied de page
        (« Imprimerie Officielle de la République Tuinsienne », avec cette faute) restée dans le texte des \term{chunks}.
  \item \textbf{Numéro de base partagé.} Le champ \texttt{article} de l'index contient le numéro de base : l'article
        220 et l'article 220 \emph{bis} partagent la valeur 220. Nous avons relevé \NAuditCollisions{} cas, dans
        \NAuditCollsources{} sources. Toute comparaison doit donc utiliser l'étiquette complète (identifiant du \term{chunk}).
\end{itemize}

\begin{table}[H]
\centering
\small
\begin{tabular}{lrrrrrrr}
\toprule
Source & Art. PDF & Art. index & PDF seul & Index seul & Sections & Échant. & Écarts \\
\midrule
ArbitrageFr & 82 & 82 & 0 & 0 & 6/6 & 6 & 0 \\
COC & 1405 & 1403 & 2 & 0 & 3/3 & 40 & 1 \\
CommerceFr & 556 & 556 & 0 & 0 & 2/2 & 12 & 3 \\
DIPFR & 76 & 75 & 1 & 0 & 2/2 & 1 & 0 \\
EnfantFr & 137 & 137 & 0 & 0 & 10/12 & 5 & 0 \\
NationaliteFr & 65 & 65 & 0 & 0 & 2/2 & 3 & 0 \\
PenalFr & 355 & 355 & 0 & 0 & 3/3 & 6 & 0 \\
Procedurecivilecom & 463 & 463 & 0 & 0 & 2/2 & 11 & 0 \\
ProcedurepenaleFr & 371 & 371 & 0 & 0 & 5/5 & 9 & 1 \\
Statutpersonnel & 222 & 222 & 0 & 0 & 16/16 & 17 & 0 \\
dtreelFr & 426 & 425 & 1 & 0 & 8/8 & 11 & 1 \\
societeFr & 493 & 489 & 4 & 0 & 6/6 & 11 & 1 \\
travailFr & 493 & 492 & 1 & 0 & 2/2 & 12 & 0 \\
tun\_RecAcFreeLi & 142 & 142 & 0 & 0 & 1/1 & 6 & 0 \\
\bottomrule
\end{tabular}
\caption{Audit des \emph{chunks} contre le texte des PDF (extraction pypdf indépendante). Étiquettes d'articles et sections (PDF/index) ; fidélité sur l'échantillon de 150 \emph{chunks}.}
\label{tab:audit}
\end{table}


\paragraph{Limite de l'extraction.} Les notes de bas de page en arabe de plusieurs PDF ne sont pas lisibles par pypdf
(glyphes non décodés) ; elles sont retirées et comptées, jamais présentées comme du texte. Les articles eux-mêmes sont en
français et l'extraction indépendante a été jugée exploitable pour les \NCorpusVerified{} sources.
